\chapter{Derive the exact probability of fixation}
\label{chap:exact}

\section{An independent answer for checking the simulations}

Repeated simulations estimate a probability. We now derive that probability
directly from the model rules. The mathematical result is exact for this
model. It is not an approximation obtained by averaging a finite number of
computer runs.

Let \(u_i\) denote the probability that A eventually fixes when the current
A count is \(i\). The subscript tells us the starting count for this question.
For example, \(u_1\) is the chance of eventual fixation from one A individual.
It is not the chance that the \emph{next event} increases the count.

The population size \(N\) and A reproductive weight \(r=1+s\) stay fixed
throughout this derivation. The transition probabilities \(T_i^+\),
\(T_i^-\), and \(T_i^0\) were derived in
Chapter~\ref{chap:moran-event}. They describe increase, decrease, and no
change during a single event.

\section{Start with the two answers we already know}

If \(i=0\), no A individuals remain. With no mutation, A cannot return.
Its probability of eventual fixation is therefore zero. If \(i=N\), all
individuals are already A, so fixation has already occurred and its
probability is one. Thus
\begin{equation}
 u_0=0,\qquad u_N=1.
 \label{eq:fixation-boundaries}
\end{equation}
These are \textbf{boundary conditions}. They give the values of the unknown
function at the two absorbing states. Our task is to find the values at
the interior states \(i=1,2,\ldots,N-1\).

\section{Why the present count is enough}

The next-event probabilities depend only on the current count and the fixed
parameters. They do not depend on the sequence of earlier counts. A random
process with this property has the \textbf{Markov property}.

Because the states are a finite list and events occur in discrete steps,
we can call this process a \textbf{finite-state Markov chain}. That name
does not add a new assumption. It names a property of the rules already
specified.

The Markov property lets us reuse the same fixation probabilities after
one event. If the next count is \(i+1\), the chance of eventual fixation
from then onward is \(u_{i+1}\). It does not matter how the process arrived
at that count.

\section{Condition on the first event}

Start at an interior state \(i\). Exactly one of three things happens next.
With probability \(T_i^+\), we move to \(i+1\). With probability
\(T_i^-\), we move to \(i-1\). With probability \(T_i^0\), we stay at
\(i\).

For each possibility, multiply the probability of that first event by the
probability of later fixation given its resulting state. Then add the three
mutually exclusive possibilities:
\begin{equation}
 u_i=T_i^+u_{i+1}+T_i^-u_{i-1}+T_i^0u_i.
 \label{eq:backward-first-step}
\end{equation}

This is called a \textbf{backward equation} or a \textbf{first-step
recurrence}. We calculate eventual success from a given starting state by
looking at the first step out of that state.

The term \(T_i^0u_i\) matters. A no-change event does not end the run.
After that event, the process has the same count, so its remaining chance
of fixation is still \(u_i\).

\begin{workedexample}
For \(N=5\), \(i=2\), and \(r=1.5\), the one-event probabilities are
\(0.3\), \(0.2\), and \(0.5\). The equation is
\[
 u_2=0.3u_3+0.2u_1+0.5u_2.
\]
There is a thirty percent chance of moving to a state whose fixation
probability is \(u_3\), a twenty percent chance of moving to one whose
fixation probability is \(u_1\), and a fifty percent chance of remaining
at the same state. None of these one-event probabilities is itself \(u_2\).
\end{workedexample}

\section{Rearrange the recurrence without skipping steps}

First subtract \(T_i^0u_i\) from both sides of
equation~\eqref{eq:backward-first-step}:
\[
 (1-T_i^0)u_i=T_i^+u_{i+1}+T_i^-u_{i-1}.
\]
Because the three transition probabilities add to one,
\(1-T_i^0=T_i^++T_i^-\). Substitute that expression:
\[
 (T_i^++T_i^-)u_i=T_i^+u_{i+1}+T_i^-u_{i-1}.
\]
Expand the left side, then move its two terms to the right:
\begin{align*}
 0&=T_i^+u_{i+1}-T_i^+u_i+T_i^-u_{i-1}-T_i^-u_i\\
  &=T_i^+(u_{i+1}-u_i)-T_i^-(u_i-u_{i-1}).
\end{align*}
Finally move the second term to the other side:
\begin{equation}
 T_i^+(u_{i+1}-u_i)=T_i^-(u_i-u_{i-1}).
 \label{eq:backward-differences}
\end{equation}
The equation now relates the improvement in fixation probability from one
extra A individual to the improvement from the preceding extra individual.

\section{Name the differences between neighbouring answers}

Let \(d_i\) denote the difference between the fixation probabilities at
counts \(i\) and \(i-1\):
\[
 d_i=u_i-u_{i-1},\qquad i=1,2,\ldots,N.
\]
For example, \(d_1=u_1-u_0=u_1\), while \(d_2=u_2-u_1\).
Equation~\eqref{eq:backward-differences} becomes
\[
 T_i^+d_{i+1}=T_i^-d_i.
\]
For an interior count, both \(i\) and \(N-i\) are positive. The
reproductive weight \(r\) is positive too, so \(T_i^+>0\). We may divide
by it:
\begin{equation}
 d_{i+1}=\frac{T_i^-}{T_i^+}d_i.
 \label{eq:difference-ratio}
\end{equation}

Now substitute the actual transition probabilities. Recall that
\(W_i=ri+N-i\) is their shared total reproductive weight:
\begin{align*}
 \frac{T_i^-}{T_i^+}
 &=\frac{\dfrac{N-i}{W_i}\dfrac{i}{N}}
          {\dfrac{ri}{W_i}\dfrac{N-i}{N}}\\[4pt]
 &=\frac{(N-i)i}{ri(N-i)}\\
 &=\frac1r.
\end{align*}
The factors \(W_i\) and \(N\) cancel first. The positive factors \(i\)
and \(N-i\) cancel next. The important simplification is that the ratio
does not depend on the current count:
\[
 d_{i+1}=\frac{d_i}{r}.
\]

\section{Follow the pattern of differences}

Starting from \(d_1\), the recurrence gives
\[
 d_2=\frac{d_1}{r},\qquad
 d_3=\frac{d_2}{r}=\frac{d_1}{r^2},\qquad
 d_4=\frac{d_1}{r^3}.
\]
Each new difference is the preceding difference multiplied by \(1/r\).
Such a pattern is a \textbf{geometric sequence}. For any allowed index \(i\),
the pattern is
\[
 d_i=d_1r^{-(i-1)}.
\]
A negative exponent means a reciprocal. For example, \(r^{-2}=1/r^2\).

To recover \(u_i\), add the first \(i\) differences. The intermediate
terms cancel:
\begin{align*}
 d_1+d_2+\cdots+d_i
 &=(u_1-u_0)+(u_2-u_1)+\cdots+(u_i-u_{i-1})\\
 &=u_i-u_0\\
 &=u_i.
\end{align*}
This cancellation is sometimes called a \textbf{telescoping sum}. In this
case it just means that every intermediate \(+u\) has an equal \(-u\).

Let \(j\) be an integer index used to count the powers in a sum. The symbol
\(\sum\) means ``add the listed terms''. We can now write
\[
 u_i=d_1\sum_{j=0}^{i-1}r^{-j}.
\]
For instance, when \(i=3\), the sum is \(1+r^{-1}+r^{-2}\).
At the other boundary, \(u_N=1\), so
\[
 1=d_1\sum_{j=0}^{N-1}r^{-j}.
\]
Dividing by the positive sum determines the unknown first difference:
\[
 d_1=\frac{1}{\displaystyle\sum_{j=0}^{N-1}r^{-j}}.
\]
Substitute this into the expression for \(u_i\):
\begin{equation}
 u_i=\frac{\displaystyle\sum_{j=0}^{i-1}r^{-j}}
           {\displaystyle\sum_{j=0}^{N-1}r^{-j}}.
 \label{eq:fixation-sum}
\end{equation}

\section{Calculate the geometric sum explicitly}

Assume for the moment that \(r\neq1\), meaning that \(r\) is not equal
to one. Let \(q=1/r\). Let \(G_i\) denote the sum of the first \(i\)
powers beginning with \(q^0=1\):
\[
 G_i=1+q+q^2+\cdots+q^{i-1}.
\]
Multiplying the whole sum by \(q\) gives
\[
 qG_i=q+q^2+\cdots+q^{i-1}+q^i.
\]
Subtract the second equation from the first. The terms in the middle cancel:
\[
 G_i-qG_i=1-q^i.
\]
Factor the left side to obtain \((1-q)G_i=1-q^i\). Since \(q\neq1\),
divide by \(1-q\):
\[
 G_i=\frac{1-q^i}{1-q}.
\]
Apply this formula to both sums in equation~\eqref{eq:fixation-sum}.
Their factors \(1-q\) cancel. Replacing \(q\) by \(1/r\) gives
\begin{equation}
 \boxed{\displaystyle u_i=\frac{1-r^{-i}}{1-r^{-N}}\quad(r\neq1).}
 \label{eq:exact-fixation}
\end{equation}
This is the exact Moran fixation probability under our stated rules.
Its single-initial-copy form is a standard result discussed, with a
log-fitness convention, by McCandlish and colleagues~\cite{mccandlish2015}.

For one initial A individual, substitute \(i=1\):
\[
 u_1=\frac{1-r^{-1}}{1-r^{-N}}.
\]
For any starting count, substitute \(r=1+s\) to express the same answer
using our advantage parameter \(s\).

\section{Handle neutrality separately}

If \(s=0\), then \(r=1\). Substituting \(r=1\) directly into
equation~\eqref{eq:exact-fixation} gives zero divided by zero, which is
undefined. That does not mean the model has no answer. It means that this
particular shortened expression needs a separate neutral case.

Return to equation~\eqref{eq:fixation-sum}. When \(r=1\), every power
\(r^{-j}\) equals one. The numerator contains \(i\) ones. The denominator
contains \(N\) ones. Therefore
\begin{equation}
 \boxed{\displaystyle u_i=\frac{i}{N}\quad(s=0).}
 \label{eq:neutral-fixation}
\end{equation}
For one initial A individual, the neutral probability is \(1/N\). The
neutral fixation probability equals the initial frequency, not one half
unless the initial frequency happens to be one half.

\section{A complete five-individual calculation}

Take \(N=5\) and \(s=0.5\), so \(r=1.5=3/2\) and \(r^{-1}=2/3\).
First compute the denominator of equation~\eqref{eq:exact-fixation}:
\[
 1-r^{-N}=1-\left(\frac23\right)^5
 =1-\frac{32}{243}=\frac{211}{243}.
\]
Starting from two A individuals, the numerator is
\[
 1-r^{-2}=1-\left(\frac23\right)^2
 =1-\frac49=\frac59.
\]
Dividing by a fraction means multiplying by its reciprocal:
\[
 u_2=\frac{5/9}{211/243}
 =\frac59\frac{243}{211}=\frac{135}{211}\approx0.6398.
\]
The symbol \(\approx\) means ``approximately equal to''. The exact fraction
is \(135/211\); the decimal is rounded.

Starting from just one A individual, the numerator is \(1-2/3=1/3\), so
\[
 u_1=\frac{1/3}{211/243}=\frac{81}{211}\approx0.3839.
\]
Even this substantial reproductive advantage does not make fixation certain.
With one initial copy, the loss probability is
\(1-u_1=130/211\approx0.6161\).

For comparison, the neutral probabilities are \(2/5=0.4\) from two copies
and \(1/5=0.2\) from one copy. Selection improves the odds in both
comparisons. The absolute probability and the improvement over neutrality
are different quantities.

\section{Check the answer using simultaneous equations}

There is another way to solve the backward equations. For a small population,
we can solve the interior equations together without using the geometric-sum
formula. This gives an independent check on the derivation and the code.

Take \(N=3\) and \(r=2\), so \(s=1\). This larger advantage is used
only to make the arithmetic short. The unknown probabilities are \(u_1\)
and \(u_2\); the boundaries are \(u_0=0\) and \(u_3=1\).

At \(i=1\), the total weight is \(2(1)+2=4\). The probabilities are
\(T_1^+=1/3\), \(T_1^-=1/6\), and \(T_1^0=1/2\). The backward
equation is
\[
 u_1=\frac13u_2+\frac16(0)+\frac12u_1.
\]
Subtract \(u_1/2\), then multiply by six, to get
\[
 3u_1=2u_2,\qquad\text{so}\qquad u_1=\frac23u_2.
\]

At \(i=2\), the total weight is \(2(2)+1=5\). The probabilities are
\(T_2^+=4/15\), \(T_2^-=2/15\), and \(T_2^0=9/15\). Hence
\[
 u_2=\frac4{15}(1)+\frac2{15}u_1+\frac9{15}u_2.
\]
Multiply by fifteen and collect the \(u_2\) terms:
\[
 15u_2=4+2u_1+9u_2,\qquad\text{so}\qquad6u_2=4+2u_1.
\]
Substitute \(u_1=2u_2/3\), then multiply by three:
\begin{align*}
 6u_2&=4+\frac43u_2,\\
 18u_2&=12+4u_2,\\
 14u_2&=12.
\end{align*}
Thus \(u_2=6/7\), and \(u_1=(2/3)(6/7)=4/7\).

The exact formula gives the same results:
\[
 u_1=\frac{1-1/2}{1-1/8}=\frac47,\qquad
 u_2=\frac{1-1/4}{1-1/8}=\frac67.
\]

For larger populations, we ask Julia to solve all \(N-1\) interior
equations together. A matrix is a rectangular arrangement of numbers that
can store the coefficients of such equations. The numerical solver handles
the elimination. We still supply the equations and the two boundary values.
Agreement with the closed formula is a useful check because the numerical
procedure does not evaluate that formula to obtain its answer.

\section{A preview of stable numerical evaluation}

Mathematically, equation~\eqref{eq:exact-fixation} is correct. A computer,
however, stores numbers with a finite number of significant digits. When
\(s\) is extremely small, both numerator and denominator subtract two
nearly equal numbers. Some useful digits can be lost.

The natural logarithm, written \(\log\), and the exponential function,
written \(\exp\), let us rewrite the powers. They are inverse functions:
\(\exp(\log r)=r\) for positive \(r\). Let
\(\sigma=\log r=\log(1+s)\). This Greek letter is read as ``sigma'' and
is the difference in log fitness used by our reference paper. It is not
exactly the same as \(s\).

Using \(r^{-i}=\exp(-i\sigma)\), the same probability is
\[
 u_i=\frac{1-\exp(-i\sigma)}{1-\exp(-N\sigma)}.
\]
Julia provides \code{log1p(s)} to calculate \(\log(1+s)\) accurately
when \(s\) is small. It provides \code{expm1(z)} to calculate
\(\exp(z)-1\) accurately when its argument, called \(z\) here, is near
zero. We can use \code{expm1} in both numerator and denominator. Both signs
reverse, so the ratio is unchanged. The neutral case still needs its explicit
answer \(i/N\).

These functions change numerical accuracy, not the biological model. The
code walkthrough explains their use after first showing the readable formula.

\begin{checkpoint}
\textbf{Question 1.} A neutral population has \(N=20\) and starts with
\(i=3\). What is the fixation probability?

\textbf{Answer.} Neutrality gives \(u_i=i/N\), so
\(u_3=3/20=0.15\). The loss probability is \(1-0.15=0.85\).

\textbf{Question 2.} In the three-individual example, verify the first
backward equation using \(u_1=4/7\) and \(u_2=6/7\).

\textbf{Answer.} Its right side is
\((1/3)(6/7)+(1/6)(0)+(1/2)(4/7)=2/7+0+2/7=4/7\), which equals its
left side \(u_1\).

\textbf{Question 3.} Why is it wrong to identify \(T_1^+\) with \(u_1\)?

\textbf{Answer.} The first quantity is the probability of an increase in
the next single event. The second is the probability of eventually reaching
the all-A boundary, possibly after many increases, decreases, and unchanged
events. An initial increase can still be followed by loss.
\end{checkpoint}

\section{What we have established}

The exact answer came from four ingredients: the event probabilities, the
Markov property, conditioning on the first event, and the two absorbing
boundaries. No simulation output was used to derive it. We can therefore
use it as an external target when assessing simulation estimates. Differences
of the expected sampling size are normal. Large systematic discrepancies
would instead suggest a mismatch in assumptions, a numerical problem, or
an implementation error.
